%% NeurIPS workshop submission template
%% Based on the official NeurIPS 2024 LaTeX style
%%
%% SETUP: download neurips_2024.sty from https://neurips.cc/Conferences/2024/PaperInformation/StyleFiles
%% or open this file on Overleaf and add neurips_2024.sty to the project.
%%
%% Compile with: pdflatex -> bibtex -> pdflatex -> pdflatex
%%
%% Common workshop page limits: 4 pages (content) + unlimited references.
%% Toggle \usepackage[final]{neurips_2024} when submitting camera-ready.

\documentclass{article}

% -----------------------------------------------------------------------
% NeurIPS style  (swap 'preprint' for 'final' on camera-ready)
% -----------------------------------------------------------------------
\usepackage[preprint]{neurips_2024}

% -----------------------------------------------------------------------
% Standard NeurIPS package stack
% -----------------------------------------------------------------------
\usepackage[utf8]{inputenc}
\usepackage[T1]{fontenc}
\usepackage{hyperref}
\usepackage{url}
\usepackage{booktabs}
\usepackage{amsfonts}
\usepackage{amsmath,amssymb}
\usepackage{nicefrac}
\usepackage{microtype}
\usepackage{xcolor}
\usepackage{graphicx}
\usepackage{natbib}

\bibliographystyle{abbrvnat}

% -----------------------------------------------------------------------
% Title
% -----------------------------------------------------------------------
\title{Dynamics of Orchestrator--Worker Large Language Model Systems}

% -----------------------------------------------------------------------
% Authors  (NeurIPS format: name, affiliation, email stacked per author)
% -----------------------------------------------------------------------
\author{%
  Jacqueline Doan \\
  Royal Bank of Canada \\
  Toronto, Ontario \\
  \texttt{jacqueline.doan@rbc.com}
  % add co-authors with \And:
  % \And
  % Author Two \\
  % Affiliation \\
  % \texttt{email@domain.com}
}

\begin{document}

\maketitle

% -----------------------------------------------------------------------
% Abstract  (NeurIPS: ~150-250 words)
% -----------------------------------------------------------------------
\begin{abstract}
% Multi-agent systems built from large language models (LLMs) route tasks
% across specialised worker agents, yet the macroscopic dynamics that emerge
% from these routing decisions remain poorly characterised.
% Here we study orchestrator--worker systems through a control-theoretic lens,
% treating the shared workspace as a state variable evolving under a switched
% dynamical system whose switching signal is the routing sequence.
% Three workers play complementary roles: a Researcher (expansion), a Critic
% (contraction), and a Synthesizer (dissipation/projection).
% We sweep routing strategies (LLM-orchestrated, random, fixed, round-robin)
% and Synthesizer placement schedules (peer, terminal, scheduled-$k$)
% across a set of multi-step reasoning tasks, recording routing sequences,
% workspace embeddings, and final answers at each step.
% Using block entropy, pre-registered motif frequencies, and inter-agent
% timing statistics, we show that LLM orchestrators self-organise into
% low-entropy, Synthesizer-avoiding limit cycles, while random routing
% approaches maximum entropy.
% Workspace-embedding divergence across replicates serves as a
% finite-time divergence-rate proxy, revealing that routing strategy
% is the dominant source of output variance.
% These results establish a quantitative framework for characterising and
% controlling the dynamical behaviour of multi-agent LLM systems.
\end{abstract}

% -----------------------------------------------------------------------
% 1. Introduction
% -----------------------------------------------------------------------
\section{Introduction}

Multi-agent LLM frameworks have become a dominant paradigm for complex
reasoning tasks, yet little is known about the macroscopic dynamics that
emerge when agents communicate through a shared workspace.
Prior work has characterised individual agent capabilities
\citep{CITE_singleagent}, prompt engineering for multi-step
reasoning \citep{CITE_chainofthought}, and benchmark performance of
collaborative agents \citep{CITE_multiagentbench}, but the dynamical
properties of the routing process itself have received less attention.

We approach this gap through the lens of switched dynamical systems
\citep{liberzon2003switching}.
Let $x(t) \in \mathbb{R}^d$ denote the workspace embedding at round $t$.
The system evolves as
\begin{equation}
  x(t+1) = f_{\sigma(t)}\!\bigl(x(t)\bigr), \qquad \sigma(t) \in \{R, C, S\},
  \label{eq:switched}
\end{equation}
where $\sigma(t)$ is the routing (switching) signal and $f_R$, $f_C$, $f_S$
are the Researcher, Critic, and Synthesizer operators respectively.

% TODO: expand with literature on agent communication, LLM stochasticity,
%       and empirical NLP dynamics

\paragraph{Contributions.}
\begin{enumerate}
  \item We formalise orchestrator--worker LLM systems as switched dynamical
        systems and identify the Synthesizer as a dissipative element.
  \item We introduce a suite of routing-sequence metrics (block entropy,
        motif frequencies, inter-arrival times) and workspace divergence
        measures applicable to any multi-agent LLM system.
  \item We show empirically that LLM orchestrators converge to low-entropy,
        Synthesizer-free limit cycles, independent of task.
\end{enumerate}

% -----------------------------------------------------------------------
% 2. Background
% -----------------------------------------------------------------------
\section{Background}

\paragraph{Switched dynamical systems.}
A switched system \citep{liberzon2003switching} is a family of subsystems
$\{f_p\}_{p \in \mathcal{P}}$ together with a switching signal
$\sigma : \mathbb{Z}_{\geq 0} \to \mathcal{P}$ that selects the active
subsystem at each step.
Stability depends jointly on the subsystem dynamics and on the switching
signal; in particular, dwell-time constraints on $\sigma$ are a standard
tool for guaranteeing stability \citep{liberzon2003switching}.

\paragraph{Control-theoretic worker roles.}
The three workers implement complementary operators on the workspace:
\begin{itemize}
  \item \textbf{Researcher} $f_R$: expansion — appends new content,
        $\|x(t+1)\| \geq \|x(t)\|$.
  \item \textbf{Critic} $f_C$: contraction — appends adversarial commentary
        opposing prior claims \citep{lohmiller1998contraction}.
  \item \textbf{Synthesizer} $f_S$: dissipation/projection —
        \emph{replaces} the workspace,
        $x(t+1) = \Pi\bigl(x(t)\bigr)$, $\|\Pi(x)\| \leq \|x\|$
        \citep{khalil2002nonlinear}.
\end{itemize}
The update rule is therefore
\begin{equation}
  x(t+1) =
  \begin{cases}
    x(t) \oplus \delta_w(t) & w \in \{R, C\} \\
    \Pi\!\bigl(x(t)\bigr)   & w = S,
  \end{cases}
  \label{eq:update}
\end{equation}
where $\oplus$ denotes text-space concatenation.

% -----------------------------------------------------------------------
% 3. Methods
% -----------------------------------------------------------------------
\section{Methods}

\paragraph{Agent architecture.}
Each worker receives a structured system prompt encoding its role and
operates on a shared plain-text workspace via the Anthropic Messages API
(\texttt{claude-haiku-4-5-20251001}).
The orchestrator implements the feedback policy
$\sigma(t) = \pi\bigl(x(t), r\bigr)$,
where $r$ is the user task (reference signal).

\paragraph{Experimental sweep.}
We sweep over routing strategies
$\in \{\text{orchestrator}, \text{random}, \text{fixed-}w, \text{round\_robin}\}$
and Synthesizer placement schedules
$\in \{\text{peer}, \text{terminal}, \text{scheduled-}k\}$.
Each cell (task $\times$ placement $\times$ strategy $\times$ $T_o$) is
replicated $k$ times with deterministic seeds.

\paragraph{Embeddings.}
Workspace and final-answer embeddings use
\texttt{all-MiniLM-L6-v2} \citep{reimers2019sentence}
($d=384$, $\ell_2$-normalised), loaded locally for reproducibility.

% -----------------------------------------------------------------------
% 4. Results
% -----------------------------------------------------------------------
\section{Results}

\subsection{Routing entropy}

The order-$n$ block entropy of the routing sequence is
\begin{equation}
  H_n = -\sum_{w \in \{R,C,S\}^n} p(w)\log_2 p(w),
  \label{eq:block_entropy}
\end{equation}
where $p(w)$ is the empirical block frequency pooled across replicates
\citep{cover2006elements}.

% TODO: report H1, H2, H3 per condition; insert figure

\subsection{Motif analysis}

We pre-registered eight routing motifs \emph{prior} to data collection
(Table~\ref{tab:motifs}) to avoid post-hoc inflation of significance.

\begin{table}[htbp]
  \centering
  \caption{Pre-registered routing motifs and observed mean counts per replicate.}
  \label{tab:motifs}
  \small
  \begin{tabular}{lll}
    \toprule
    Motif & Interpretation & Orchestrator / Random \\
    \midrule
    RC     & Researcher--Critic alternation, period 2 & \\
    RCRC   & Period-2 limit cycle                     & \\
    RCSRCS & Three-phase rhythm (conjectured healthy)  & \\
    RSRS   & Premature consolidation                  & \\
    SS     & Synthesizer doublet                      & \\
    SSSS   & Degenerate self-feeding                  & \\
    RR     & Researcher monoculture                   & \\
    CC     & Critic monoculture                       & \\
    \bottomrule
  \end{tabular}
\end{table}

\subsection{Workspace divergence}

For $K$ replicates, the finite-time divergence-rate proxy at round $t$ is
\begin{equation}
  \Delta(t) = \frac{1}{\binom{K}{2}}
              \sum_{i < j} \bigl(1 - \cos(x_i(t),\, x_j(t))\bigr).
  \label{eq:divergence}
\end{equation}
Increasing $\Delta(t)$ indicates sensitive dependence on the random seed,
analogous to the finite-time Lyapunov exponent \citep{strogatz2015nonlinear}.

% TODO: insert divergence-by-round figure

% -----------------------------------------------------------------------
% 5. Discussion
% -----------------------------------------------------------------------
\section{Discussion}

% TODO

% -----------------------------------------------------------------------
% 6. Conclusion
% -----------------------------------------------------------------------
\section{Conclusion}

% TODO

% -----------------------------------------------------------------------
% Acknowledgements  (unnumbered per NeurIPS style)
% -----------------------------------------------------------------------
\section*{Acknowledgements}

% TODO

% -----------------------------------------------------------------------
% References
% -----------------------------------------------------------------------
\bibliography{references}

\end{document}
